% Part: first-order-logic
% Chapter: model-theory
% Section: nonstandard-arithmetic

\documentclass[../../../include/open-logic-section]{subfiles}

\begin{document}

\olfileid[hi]{mod}{bas}{nsa}
\section{अंकगणित के गैर-मानक निदर्श}

\begin{defn}
मान लें $\Lang{L_N}$ अंकगणित की वह !!{language} है जिसमें एक !!{constant}
$\Obj{0}$, एक द्विस्थानी !!{predicate} $<$, एक एकस्थानी !!{function}
$\prime$, और द्विस्थानी !!{function}s $+$ तथा $\times$ हैं।
\begin{enumerate}
\item अंकगणित का \emph{मानक निदर्श}, $\Lang{L_N}$ की वह !!{structure}
  $\Struct{N}$ है जिसमें $\Nat = \{ 0, 1, 2, \dots\}$ है और जो $\Obj{0}$
  का निर्वचन $0$, $<$ का~$\Nat$ पर लघुतर संबंध, तथा $\prime$, $+$ और
  $\times$ का निर्वचन क्रमशः $\Nat$ पर उत्तरवर्ती, योग तथा गुणन के रूप में
  करती है।
\item \emph{सत्य अंकगणित}, सिद्धांत $\Theory{N}$ है।
\end{enumerate}
\end{defn}

$\Lang{L_N}$ में काम करते समय हम $\Obj{0}$ पर उत्तरवर्ती फलन के $n$
अनुप्रयोगों वाले रूप $\Obj{0}^{\prime\dots\prime}$ के प्रत्येक पद को
संक्षेप में~$\num{n}$ लिखते हैं।

\begin{defn}
$\Lang{L_N}$ की कोई !!{structure}~$\Struct{M}$ \emph{मानक} तभी और केवल तभी
है जब $\Struct{N} \iso \Struct{M}$।
\end{defn}

\begin{thm}\ollabel{thm:non-std}
सत्य अंकगणित के गैर-मानक !!{enumerable} निदर्श होते हैं।
\end{thm}

\begin{proof}
नया !!{constant}~$c$ प्रस्तुत करके $\Lang{L_N}$ का विस्तार करें और सिद्धांत
\[
\Theory{N} \cup \Setabs{\num{n} < c}{n \in \Nat}.
\]
पर विचार करें। यह सिद्धांत परिमित रूप से संतुष्टनीय है, अतः संहतता से इसका
एक निदर्श $\Struct{M}$ है, जिसे अधोमुखी लोवेनहाइम--स्कोलेम प्रमेय से
!!{enumerable} लिया जा सकता है। जहाँ $\Domain{M}$, $\Struct{M}$ का प्रांत
है, $\Struct{M}$ को $\Lang{L}$ के अतार्किक अचरों का निर्वचन इस प्रकार करने
दें: $\mathbf{z} = \Assign{\Obj{0}}{M} \in \Domain{M}$, ${\prec} =
\Assign{<}{M} \subseteq M^2$, $* = \Assign{\prime}{M} \colon \Domain{M}
\to \Domain{M}$, और $\oplus = \Assign{+}{M}, \otimes =
\Assign{\times}{M} \colon \Domain{M}^2 \to \Domain{M}$। प्रत्येक $x \in
\Domain{M}$ के लिए, $*$ को $x$ पर लागू करने से मिले $\Domain{M}$ के अवयव को
हम $x^*$ लिखते हैं।

अब यदि $h$, $\Struct{N}$ और $\Struct{M}$ के बीच समाकृतिकता होता, तो कोई $n
\in \Nat$ होता ऐसा कि $h(n) = \Assign{c}{M}$। अतः $\Struct{N}$ में कोई
ऐसा नियतन $s$ लें कि $s(x) = n$। तब $\Sat{N}{\eq[\num{n}][x]}[s]$;
\olref[iso]{thm:isom} के प्रमाण से $\Sat{M}{\eq[\num{n}][x]}[h \circ s]$
भी, अतः $\Assign{c}{M} = \mathbf{z}^{*\cdots *}$ ($*$ को $n$ बार दोहराकर)।
किंतु यह असंभव है, क्योंकि मान्यता से $\Sat{M}{\num{n} < c}$ और $\prec$
अस्वपरावर्ती है। इसलिए $\Struct{M}$ गैर-मानक है।
\end{proof}

\begin{prob}
किसी समुच्चय $X$ पर संबंध $R$ \emph{सुस्थापित} तभी और केवल तभी है जब~$R$
में कोई अनंत अवरोही शृंखला न हो, अर्थात् $X$ में ऐसे $x_0$, $x_1$,
$x_2$,~\dots न हों कि $\dots x_2Rx_1Rx_0$। ज़र्मेलो--फ़्रेंकेल समुच्चय
सिद्धांत~$ZF$ को संगत मानते हुए दिखाइए कि $ZF$ के गैर-सुस्थापित निदर्श होते
हैं, अर्थात् ऐसे निदर्श $\mathfrak{M}$ कि $\dots x_2 \in x_1 \in x_0$।
\end{prob}

चूँकि \olref{thm:non-std} का गैर-मानक निदर्श $\Struct{M}$ मानक निदर्श से
प्राथमिकतः तुल्य है, $\Struct{M}$ के अनेक गुण निकाले जा सकते हैं। इस खंड का
शेष भाग इसी कार्य को समर्पित है, जिससे हम $\Theory{N}$ के !!{enumerable}
गैर-मानक निदर्शों का सटीक अभिलक्षण प्राप्त कर सकेंगे।

\begin{enumerate}
\item $\Domain{M}$ का कोई सदस्य स्वयं से $\prec$-लघुतर नहीं है: वाक्य
  $\lforall[x][\lnot x < x]$, $\Struct{N}$ में और इसलिए $\Struct{M}$ में
  सत्य है।
\item इसी प्रकार के तर्क से हमें मिलता है कि $\prec$, $\Domain{M}$ का एक
  \emph{रैखिक क्रम} है, अर्थात् $\Domain{M}$ पर संपूर्ण, अस्वपरावर्ती,
  संक्रामी संबंध।
\item अवयव $\mathbf{z}$, $\Domain{M}$ का $\prec$-लघुतम अवयव है।
\item $\Domain{M}$ का प्रत्येक सदस्य अपने $*$-उत्तरवर्ती से $\prec$-लघुतर
  है और $x^*$, $x$ से बड़ा $\Domain{M}$ का $\prec$-लघुतम सदस्य है।
\item $\Struct{M}$ में $\Nat$ से समाकृतिक ($\prec$ का) एक आरंभिक खंड है:
  $\mathbf{z}, \mathbf{z}^*, \mathbf{z}^{**}, \dots$, जिसे हम
  $\Domain{M}$ का \emph{मानक भाग} कहते हैं। $\Domain{M}$ का कोई अन्य सदस्य
  \emph{गैर-मानक} है। $\Domain{M}$ में गैर-मानक सदस्य अवश्य होने चाहिए,
  अन्यथा \olref{thm:non-std} के प्रमाण का फलन $h$ एक समाकृतिकता है। हम $n,
  m, \dots$ को $\Struct{M}$ के इस मानक भाग पर विचरने वाले !!{variable}s के
  रूप में प्रयोग करते हैं।
\item प्रत्येक गैर-मानक अवयव किसी भी मानक अवयव से बड़ा है; इसका कारण है कि
  प्रत्येक $n \in \Nat$ के लिए,
  \[
  \Sat{N}{\lforall[z][(\lnot(\eq[z][\Obj{0}] \lor
  \dots \lor \eq[z][\num{n}]) \lif \num{n} < z)]},
  \]
  अतः यदि $z \in \Domain{M}$ सभी मानक अवयवों से भिन्न है, तो वह उन सभी से
  \emph{बड़ा} होना चाहिए।
\item $\mathbf{z}$ के अतिरिक्त $\Domain{M}$ का प्रत्येक सदस्य,
  $\Domain{M}$ के किसी अद्वितीय अवयव का $*$-उत्तरवर्ती है, जिसे $^*x$ से
  निरूपित करते हैं। यदि $x = y^*$, तो $x$ और $y$ में से एक के मानक होने पर
  दोनों मानक हैं (और एक के गैर-मानक होने पर दोनों गैर-मानक)।
\item $\Domain{M}$ पर तुल्यता संबंध $\approx$ को यह कहकर परिभाषित करें कि
  $x \approx y$ तभी और केवल तभी जब किसी \emph{मानक} $n$ के लिए या तो $x
  \oplus n = y$ या $y \oplus n =x$। दूसरे शब्दों में $x \approx y$ तभी और
  केवल तभी जब $x$ और $y$ परिमित दूरी पर हों। यदि $n$ और $m$ मानक हैं, तो
  $n \approx m$। $x$ के \emph{खंड} को तुल्यता वर्ग $[x] = \Setabs{y}{x
  \approx y}$ परिभाषित करें।
\item मान लें $x \prec y$, जहाँ $x \not\approx y$। चूँकि
  $\Sat{N}{\lforall[x][\lforall[y][(x < y \lif (x' < y \lor x' =
      y))]]}$, इसलिए या तो $x^* \prec y$ या $x^* = y$। दूसरा असंभव है,
  क्योंकि उससे $x \approx y$ निष्पन्न होता है, अतः $x \prec y$। इसी प्रकार
  यदि $x \prec y$ और $x \not\approx y$, तो $x \prec {^*y}$। इसलिए यदि $x
  \prec y$ और $x \not\approx y$, तो प्रत्येक $w \approx x$, प्रत्येक $v
  \approx y$ से $\prec$-लघुतर है। तदनुसार प्रत्येक खंड $[x]$ द्वि-अनंत
  शृंखला बनाता है
  \[
  \cdots \prec  {^{**}x} \prec {^*}x \prec x \prec x^* \prec x^{**}
  \prec \cdots
  \]
  जिसे $Z$-शृंखला कहते हैं क्योंकि इसका क्रम-प्रकार पूर्णांकों वाला है।
\item $\prec$ क्रम को खंडों तक उठाया जा सकता है: यदि $x \prec y$, तो $x$ का
  खंड $y$ के खंड से लघुतर है। कोई खंड \emph{गैर-मानक} है यदि उसमें कोई
  गैर-मानक अवयव हो। मानक खंड सबसे छोटा खंड है।
\item कोई लघुतम गैर-मानक खंड नहीं है: यदि $y$ गैर-मानक है, तो कोई $x \prec
  y$ है जहाँ $x$ भी गैर-मानक और $x \not\approx y$ है। प्रमाण: मानक निदर्श
  $\Struct{N}$ में प्रत्येक संख्या दो से विभाज्य है, संभवतः शेषफल एक के साथ,
  अर्थात्
  $\Sat{N}{\lforall[y][\lforall[x][(y = x + x \lor y = x + x +
        \Obj{0}')]]}$। प्राथमिक तुल्यता से प्रत्येक $y \in \Domain{M}$ के
  लिए कोई $x \in \Domain{M}$ है ऐसा कि या तो $x \oplus x = y$ या $x
  \oplus x \oplus \mathbf{z}^*= y$। यदि $x$ मानक होता, तो $y$ भी मानक
  होता; अतः $x$ गैर-मानक है। इसके अतिरिक्त $x$ और $y$ भिन्न खंडों में हैं,
  अर्थात् $x \not\approx y$। इसे देखने के लिए मान लें कि वे एक ही खंड में
  हैं, अर्थात् किसी मानक $n$ के लिए $x \oplus n = y$। यदि $y = x \oplus
  x$, तो $x \oplus n = x \oplus x$, जिससे योग के निरस्तीकरण नियम द्वारा $x
  = n$ (यह नियम $\Struct{N}$ में और इसलिए $\Struct{M}$ में भी सत्य है), और
  अंततः $x$ मानक होता। $y = x \oplus x \oplus \mathbf{z}^*$ होने पर भी इसी
  प्रकार।
\item इसी प्रकार के तर्क से कोई महत्तम खंड नहीं है।
\item खंडों का क्रम घना है: यदि $[x]$, $[y]$ से लघुतर है (जहाँ $x
  \not\approx y$), तो उन दोनों से भिन्न कोई खंड $[z]$ उनके बीच है। मान लें
  $x \prec y$। पहले की तरह $x \oplus y$ दो से विभाज्य है (संभवतः शेषफल के
  साथ), अतः कोई $u \in \Domain{M}$ है ऐसा कि या तो $x \oplus y = u \oplus
  u$ या $x \oplus y = u \oplus u \oplus \mathbf{z}^*$। अवयव $u$, $x$ और
  $y$ का औसत है, अतः उनके बीच है। मान लें $x \oplus y = u \oplus u$ (दूसरी
  स्थिति भी इसी प्रकार): यदि $u \approx x$, तो किसी मानक $n$ के लिए:
  \[
  x \oplus y = x \oplus n \oplus x \oplus n,
  \]
  अतः $y = x \oplus n \oplus n$ और हमें $x \approx y$ मिलता, जो मान्यता के
  विरुद्ध है। हम निष्कर्ष निकालते हैं कि $u \not\approx x$। इसी प्रकार का
  तर्क $u \not\approx y$ देता है।
\end{enumerate}
अतः गैर-मानक खंड परिमेय संख्याओं की तरह क्रमित हैं: वे अंतबिंदु-रहित
!!a{enumerable} रैखिक क्रम बनाते हैं। इससे निष्पन्न होता है
कि सत्य अंकगणित के किन्हीं दो !!{enumerable} गैर-मानक निदर्शों $\Struct{M}_1$ और
$\Struct{M_2}$ के केवल $<$ तथा $=$ वाली भाषा में अपचय समाकृतिक हैं। वास्तव
में समाकृतिकता $h$ इस प्रकार परिभाषित की जा सकती है: $\Struct{M_1}$ और
$\Struct{M_2}$ के मानक भाग मानक निदर्श $\Struct{N}$ से और इसलिए परस्पर
समाकृतिक हैं। गैर-मानक भाग बनाने वाले खंड स्वयं
परिमेय संख्याओं की तरह क्रमित हैं और इसलिए \olref[dlo]{thm:cantorQ} से
समाकृतिक हैं; प्रत्येक में मनमाने अवयवों का मेल करके और फिर $\Struct{M_1}$
में $x$ के उत्तरवर्ती की छवि को $\Struct{M_2}$ में $x$ की छवि का उत्तरवर्ती
लेकर, खंडों की समाकृतिकता को खंडों के \emph{भीतर} समाकृतिकता तक विस्तारित
किया जा सकता है। ध्यान दें कि इससे यह निष्पन्न \emph{नहीं} होता कि
$\mathfrak{M}_1$ और $\mathfrak{M}_2$ अंकगणित की पूर्ण भाषा में समाकृतिक हैं
(वास्तव में समाकृतिकता सदा किसी चिह्नक-समूह के सापेक्ष होती है), क्योंकि
$\Domain{M_1}$ और $\Domain{M_2}$ पर योग और गुणन परिभाषित करने के
गैर-समाकृतिक ढंग होते हैं। (यह वॉट के एक प्रसिद्ध प्रमेय से भी निष्पन्न होता
है कि किसी पूर्ण सिद्धांत के गणनीय निदर्शों की संख्या~2 नहीं हो सकती।)

\begin{prob}
दिखाइए कि अंकगणित के किसी गैर-मानक निदर्श में कोई महत्तम खंड नहीं हो सकता।
\end{prob}

\begin{prob} 
मान लें $\Lang{L}$ वह प्रथम-क्रम !!{language} है जिसमें $\eq$ के अतिरिक्त
$<$ ही एकमात्र !!{predicate} है, और मान लें $\Struct{N} = (\Nat, <)$।
$\Struct{N}$ के सभी परिमित या सहपरिमित उपसमुच्चय परिभाष्य हैं। दिखाइए कि ये
$\Struct{N}$ के \emph{एकमात्र} परिभाष्य उपसमुच्चय हैं।
 
(संकेत: पहले $\Obj{prc}(x,y)$ को वह $\Lang{L}$-सूत्र मानें जो ``$x$, $y$
का ठीक पूर्ववर्ती है:'' को संक्षेपित करता है
\[
x<y \land \lnot \lexists[z][(x<z \land z < y)].
\]
अब $\Struct{N}$ के प्रत्येक परिभाष्य उपसमुच्चय के अनुरूप $\Lang{L}$ में कोई
सूत्र $!A(x)$ है। किसी भी ऐसे $!A$ के लिए वाक्य $!D$ पर विचार करें:
\[
\lexists[x][\lforall[y][\lforall[z][((x<y \land x<z \land \Obj{prc}(y,z)
  \land !A(y)) \lif !A(z))]]].
\]
दिखाइए कि $\Sat{N}{!D}$ तभी और केवल तभी जब $!A$ द्वारा परिभाषित
$\Struct{N}$ का उपसमुच्चय या तो परिमित या सहपरिमित हो।

अब $\Struct{M}$ को $\Struct{N}$ से प्राथमिकतः तुल्य कोई गैर-मानक निदर्श
मानें। यदि $a \in \Domain{M}$ गैर-मानक है, तो $b, c \in \Domain{M}$ को
$a$ से बड़ा मानें और $b$ को~$c$ का ठीक पूर्ववर्ती मानें। तब $\Domain{M}$
का कोई स्वसमाकृतिकता $h$ है ऐसा कि $h(b)=c$ (क्यों?)। अतः यदि $b$, $!A$ को
संतुष्ट करता है, तो $c$ भी करता है (क्यों?)। इससे निष्पन्न होता है कि $!D$,
$\Struct{M}$ में और इसलिए $\Struct{N}$ में भी सत्य है। किंतु इससे निष्पन्न
होता है कि $!A$ द्वारा परिभाषित $\Struct{N}$ का उपसमुच्चय या तो परिमित या
सहपरिमित है।
\end{prob}

\end{document}
